\documentclass[10pt,a4paper]{amsart}
\usepackage[margin=19mm]{geometry}
\usepackage{amsmath,amssymb,mathtools}
\usepackage[scr=rsfs,cal=euler]{mathalfa}
\usepackage{xfrac}
\usepackage[unicode,hidelinks,bookmarks=false]{hyperref}
\newcommand{\ie}{i.e.\ }
\newcommand{\cf}{cf.\ }
\newcommand{\resp}{respectively\ }
\newcommand{\Subsection}{\subsection}
\providecommand{\nobreakdash}{\nobreak\hbox{-}}
\setlength{\emergencystretch}{2em}
\multlinegap0pt
\usepackage{amssymb}
\usepackage{enumerate}
\newtheorem{lemma}{Lemma}
\newcommand{\Exp}{\mathbb{E}}
\newcommand{\Prb}{\mathbb{P}}
\newcommand{\bh}{\mathbf{H}}
\title{An entropic analogue of the MMS conjecture}
\author{Jianhang Ai, Ond\v{r}ej Ku\v{z}elka, Christos Pelekis}
\date{Source: arXiv:2606.30486v1 (CC BY 4.0)}
\begin{document}
\maketitle

\noindent\emph{Source and licence.} An entropic analogue of the MMS conjecture. Version arXiv:2606.30486v1 (CC BY 4.0), distributed by arXiv under the Creative Commons Attribution 4.0 International licence, http://creativecommons.org/licenses/by/4.0/, as recorded in the arXiv licence metadata for this exact version and shown on the abstract page https://arxiv.org/abs/2606.30486v1. Full licence notice: \texttt{licenses/arxiv-2606.30486-CC-BY-4.0.txt}.\par\smallskip

\noindent\textit{Editorial scope.} Counted unit: the article's lemma last\_lemma (source0.tex lines 460--467). The article's complete original proof of it is delivered with it (source0.tex lines 468--555, one \texttt{proof} environment of 88 lines). No mathematics is written, repaired or re-derived: the statement and the proof are the article's own lines, in the article's own notation, and no retained line is substituted or re-flowed.  No further source-local context is needed: every cross-reference of every retained line resolves inside the two delivered spans.  Nothing was omitted from any retained span, so the package carries no omission record.  No retained line contains a citation, so the wrapper carries no bibliography and no bibliography entry.  No retained line contains a figure environment, an image-inclusion command or drawing code, so no image is delivered and the package declares no asset dependency.  Editorial formatting only, in addition: an AMS \texttt{amsart} preamble that restates the article's own declarations and macros used by the retained lines, namely the environments \texttt{lemma} on separate counters, together with the standard packages those lines need.  The retained environments print in this document as Lemma 1 in order of appearance, whereas the article prints the same items with the numbering of its own full document.  The complete native source of the article as distributed in the arXiv e-print for this exact version is bundled unchanged as \texttt{sources/204024/source0.tex}.
\begin{lemma}\label{last_lemma}
Let $n\ge 2$ and $k\in [n-1]$. Then for all $i\in [n-1]$ it holds 
\[
\bh(H_i) \ge \bh(\text{Ber}(k/n)) \, ,
\]
where $H_i\sim\text{Hyp}(n,i,k)$ and 
$\text{Ber}(k/n)$ is a Bernoulli random variable of mean $k/n$. 
\end{lemma}

\begin{proof} 
The result is clearly true when $n=2$, so we may assume that $n\ge 3$. 
Note that the probability mass function of $H_1$ is equal to $(\frac{k}{n}, 1 - \frac{k}{n})$, and therefore it holds 
$\bh(H_1)=\bh(\text{Ber}(k/n))$. 
Note also that  for each $i\in [n-1]$ it holds 
$k-H_i \sim\text{Hyp}(n,n-i,k)$. Hence $k-H_i \sim H_{n-i}$ and so we have $\bh(H_i) = \bh(H_{n-i})$. 
In particular, we may assume that $i\le \lfloor n/2\rfloor$, and it  is therefore  enough to show that 
\begin{equation}\label{enough1}
\bh(H_{i+1}) \ge \bh(H_i) \, , \, \text{ for } \, i \le \lfloor n/2\rfloor -1 \, .
\end{equation}
We prove~\eqref{enough1} via a coupling argument.  
More concretely, let $\mathbb{I}\in\binom{[n]}{k}$ be chosen uniformly at random and, for each $i\in [n-1]$, consider the random variable  
\[
Z_i = |\mathbb{I} \cap [i]| \, .
\]   
Note that $H_i \sim Z_i$, for all $i\in [n-1]$, and therefore it is enough to show that 
\begin{equation}\label{enough1.2}
\bh(Z_{i+1}) \ge \bh(Z_i) \, , \, \text{ for } \, i \le \lfloor n/2\rfloor -1 \, .
\end{equation}
Observe that it holds 
\[
Z_{i+1} = Z_i + \mathbf{1}_{\{i+1\in \mathbb{I}\}} \, ,
\]
where $\mathbf{1}_{\{i+1\in \mathbb{I}\}}$ denotes the indicator of the event $\{i+1\in \mathbb{I}\}$. 
In other words, if we know the pair $(Z_i, \mathbf{1}_{\{i+1\in \mathbb{I}\}})$ then we can recover $Z_{i+1}$, and if we know the pair 
$(Z_{i+1}, \mathbf{1}_{\{i+1\in \mathbb{I}\}})$ then we can recover $Z_i$. 
This implies that 
\[
\bh(Z_i, \mathbf{1}_{\{i+1\in \mathbb{I}\}}) = \bh(Z_{i+1}, \mathbf{1}_{\{i+1\in \mathbb{I}\}}) 
\]
which, using the chain rule, can be equivalently written as 
\[
\bh(Z_i) + \bh(\mathbf{1}_{\{i+1\in \mathbb{I}\}}\mid Z_i) = 
\bh(Z_{i+1}) + \bh(\mathbf{1}_{\{i+1\in \mathbb{I}\}}\mid Z_{i+1}) \, .
\]
Hence, in order to prove~\eqref{enough1.2}, it is enough to show that 
\begin{equation}\label{enough2}
\bh(\mathbf{1}_{\{i+1\in \mathbb{I}\}}\mid Z_i) \ge \bh(\mathbf{1}_{\{i+1\in \mathbb{I}\}}\mid Z_{i+1}) \, , \, \text{ for } \, i \le \lfloor n/2\rfloor -1 \, .
\end{equation}
Let $i\le \lfloor n/2\rfloor -1$ be a positive integer. Note that 
\[
\Prb(\mathbf{1}_{\{i+1\in\mathbb{I}\}}=1 \mid Z_i = j)= \frac{\binom{i}{j}\binom{n-i-1}{k-j-1}/\binom{n}{k}}{\binom{i}{j}\binom{n-i}{k-j}/\binom{n}{k} }
= \frac{k-j}{n-i} \, .
\]
Hence, if  $h(\cdot)$ denotes the binary entropy function, we deduce that
\begin{eqnarray*}
\bh(\mathbf{1}_{\{i+1\in \mathbb{I}\}}\mid Z_i) &=& \sum_{j} \bh(\mathbf{1}_{\{i+1\in \mathbb{I}\}}\mid Z_i=j) \cdot \Prb(H_i=j)\\
&=& \sum_{j} h\left( \frac{k-j}{n-i}\right) \cdot \Prb(Z_i=j) \\
&=& \Exp\left(h\left(\frac{k-Z_i}{n-i} \right) \right) \\
&=& \Exp\left(h\left(\frac{Z_{n-i}}{n-i} \right) \right) \, ,
\end{eqnarray*}
where the last inequality follows from the fact that $k-Z_i\sim k-H_i\sim H_{n-i}\sim Z_{n-i}$. 
Similarly, we have 
\[
\Prb(\mathbf{1}_{\{i+1\in\mathbb{I}\}}=1 \mid Z_{i+1} = \ell)= \frac{\binom{i}{\ell-1}\binom{n-i-1}{k-\ell}/\binom{n}{k}}{\binom{i+1}{\ell}\binom{n-i-1}{k-\ell}/\binom{n}{k} }
= \frac{\ell}{i+1} \, , 
\]
and therefore we may write 
\[
\bh(\mathbf{1}_{\{i+1\in \mathbb{I}\}}\mid Z_{i+1}) = \Exp\left(h\left(\frac{Z_{i+1}}{i+1}\right)\right) \, .
\]
Summarizing the above, it follows from~\eqref{enough2} that it is enough to show that 
\begin{equation}\label{enough3}
\Exp\left(h\left(\frac{Z_{n-i}}{n-i} \right) \right) \ge \Exp\left(h\left(\frac{Z_{i+1}}{i+1}\right)\right) \, , \, \text{ for } \, i \le \lfloor n/2\rfloor -1 \, .
\end{equation}
We prove a bit more. Namely, we show that the sequence 
\[
E_m := \Exp\left(h\left(\frac{Z_{m}}{m}\right)\right) \, , \text{ for } \, m=1,\ldots,n-1
\]
is non-decreasing. The desired inequality~\eqref{enough3} then follows upon observing that the assumption $i \le \lfloor n/2\rfloor -1$ implies that $n-i\ge i+1$. To this end, let $m\in [n-2]$ and note that, since $Z_{m+1} = Z_m + \mathbf{1}_{\{m+1\in \mathbb{I}\}}$, we may write 
\begin{eqnarray*}
\Exp\left(\frac{Z_m}{m} \mid Z_{m+1}=\ell \right) &=&  
\Exp\left(\frac{Z_{m+1}-\mathbf{1}_{\{m+1\in\mathbb{I}\}}}{m} \mid Z_{m+1}=\ell \right)\\
&=& \frac{\ell}{m} - \frac{1}{m}\cdot \Prb(\mathbf{1}_{\{m+1\in\mathbb{I}\}}= 1 \mid Z_{m+1}=\ell)\\
&=& \frac{\ell}{m} -\frac{1}{m}\cdot \frac{\ell}{m+1}\\
&=& \frac{\ell}{m+1} \, .
\end{eqnarray*}
Since the binary entropy function is concave, the conditional  Jensen inequality yields 
\[
\Exp\left(h\left(\frac{Z_m}{m}\right) \mid Z_{m+1}=\ell \right) \,\le\, h\left(\Exp\left(\frac{Z_m}{m} \mid Z_{m+1}=\ell\right)\right) = h\left( \frac{\ell}{m+1}\right) \, .
\]
In particular, this implies that 
\[
E_m=\Exp\left(h\left(\frac{Z_m}{m}\right)\right) = 
\Exp\left( \Exp\left(h\left(\frac{Z_m}{m}\right) \mid Z_{m+1}  \right)\right) \le \Exp\left(h\left(\frac{Z_{m+1}}{m+1}\right)  \right)= E_{m+1} \, ,
\]
as desired.  
\end{proof}

\end{document}
